\begin{frame}{자주 하는 실수 (17.2)}
  \begin{itemize}
    \item \textbf{``프롬프트 = 질문 한 줄''로 쓰기} --- 역할
      지정 없으면 모델이 ``어떤 식의 번역인가''를 추측.
      프롬프트는 \textbf{프로그램}: 역할·작업·형식·예시를
      명시할수록 출력이 예측 가능.
    \item \textbf{few-shot 예시의 형식을 대충} --- 3개 중
      형식이 섞이면 \textbf{어떤} 형식을 따라야 할지 혼란.
      형식은 \textbf{모두 동일}해야.
    \item \textbf{``프롬프트가 길수록 좋다''} --- 컨텍스트에
      한계($n^2$ 계산량)가 있고, 지나치게 긴 프롬프트에서는
      \textbf{중간 부분} 정보의 영향이 오히려 줄어든다
      (\kb{lost in the middle}, Liu et al., 2023). 핵심
      지시문은 \textbf{앞이나 뒤}에 두는 것이 안전.
    \item \textbf{``temperature 0이면 환각이 사라진다''} ---
      $T=0$은 \textbf{결정적}을 의미할 뿐, \textbf{정확}은
      아님. 틀린 사실을 \textbf{확신 있게, 매번 같은
      방식으로 반복}하는 것이 $T=0$의 환각.
    \item \textbf{``단계별로 생각해봐''를 모든 작업에
      붙이기} --- ``오늘 날씨는?'' 같은 단순 사실 조회에
      CoT를 붙이면 \textbf{답은 그대로인데} 토큰 수(= 비용
     과 지연)만 늘어남. \kb{단계가 있는} 문제(산술, 논리,
      계획)에서만 효과.
  \end{itemize}
\end{frame}
